\chapter{PENDAHULUAN}
\section{Latar Belakang Masalah}

\vspace{0.5cm}
\noindent\fbox{\begin{minipage}{\textwidth}
\textit{\textbf{Latar Belakang Masalah WAJIB terdiri dari 4 komponen:}}
\begin{enumerate}[label=\alph*)]
    \item Permasalahan umum (contoh: masalah industri secara luas)
    \item Permasalahan khusus (contoh: aspek spesifik dari masalah umum)
    \item Artikel ilmiah \& teori terkait --- MINIMAL 3 artikel ilmiah
    \item Solusi yang ditawarkan oleh aplikasi StartUp yang dibuat
\end{enumerate}
\end{minipage}}
\vspace{0.5cm}

\section{Identifikasi Permasalahan}

\vspace{0.5cm}
\noindent\fbox{\begin{minipage}{\textwidth}
\textit{\textbf{Pedoman Penulisan:} Garis besar alasan pembuatan StartUp dalam penyelesaian permasalahan.}
\end{minipage}}
\vspace{0.5cm}

\section{Perumusan Masalah}

\vspace{0.5cm}
\noindent\fbox{\begin{minipage}{\textwidth}
\textit{\textbf{Pedoman Penulisan:} Berbentuk kalimat tanya (boleh lebih dari 1). Contoh: "Bagaimana mengurangi rantai distribusi hasil pertanian dengan efisien?"}
\end{minipage}}
\vspace{0.5cm}

\section{Tujuan dan Manfaat}

\vspace{0.5cm}
\noindent\fbox{\begin{minipage}{\textwidth}
\textit{\textbf{Pedoman Penulisan:} Menjelaskan tujuan dan manfaat (untuk penulis, pengguna, dan pembaca).}
\end{minipage}}
\vspace{0.5cm}

\section{Metode Penelitian}
\subsection{Teknik Pengumpulan Data}
\subsubsection*{A. Observasi}
\subsubsection*{B. Wawancara}
\subsubsection*{C. Studi Pustaka}
\subsection{Model Pengembangan Aplikasi}

\vspace{0.5cm}
\noindent\fbox{\begin{minipage}{\textwidth}
\textit{\textbf{Pedoman Penulisan:} Sebagai dasar biasanya menggunakan RUP (Rational Unified Process) atau lainnya.}
\end{minipage}}
\vspace{0.5cm}

\section{Ruang Lingkup}
